\chapter{Conclusiones}

La presente investigación ha logrado modelar matemáticamente el sistema de transporte de la Zona Metropolitana de Guadalajara (ZMG) mediante una aproximación dual, diferenciando entre la infraestructura física (Espacio $L$) y la red de servicio proyectada (Espacio $P$). Este enfoque permitió evaluar una proposición principal única: \textbf{la ZMG combina una base física sensible a fallas con una representación de servicio relativamente robusta bajo el modelo \(P\), efecto explicado por la redundancia entre rutas.}

\section{Síntesis de Hallazgos y Verificación de la Tesis Central}

El análisis topológico realizado arroja tres conclusiones principales que caracterizan el estado actual del sistema de transporte en la ZMG:

\begin{enumerate}
    \item \textbf{Topología del Espacio $L$ (Infraestructura):} El análisis confirma que la infraestructura subyacente posee una estructura cuasi-arbórea (con un coeficiente de clustering bajo, \(\approx 0.1088\)) y una conectividad frágil. Matemáticamente, la red depende críticamente de nodos articuladores centrales y corredores específicos. La baja redundancia de ciclos en este espacio implica que la infraestructura física carece de plasticidad ante fallos locales, resultando en una degradación rápida de la eficiencia global, de aproximadamente \(18.50\%\), ante ataques dirigidos a los nodos de mayor centralidad. En contraste, los fallos aleatorios retienen en promedio \(97.17\%\) de la eficiencia tras 50 remociones, por lo que la fragilidad principal se asocia a nodos estructuralmente críticos y no a cualquier pérdida puntual.

    \item \textbf{Topología del Espacio $P$ (Servicio):} En contraste, la proyección al Espacio $P$ revela una red de Mundo Pequeño (con un diámetro $D=6$) y una alta densidad de \textit{cliques} locales. Esta estructura mitiga las deficiencias del Espacio $L$ dentro del modelo: la superposición de múltiples rutas sobre los mismos tramos físicos genera redundancia estructural que mantiene la conectividad del componente gigante al 98.74\% incluso tras la remoción de los 50 supernodos más importantes.

    \item \textbf{Verificación de la Tesis Central:} Existe una clara desconexión entre la eficiencia topológica de la infraestructura (baja) y la de la proyección de servicio (alta). En consecuencia, la evidencia respalda la proposición central del trabajo en el marco modelado: \textbf{la ZMG combina una base física sensible a fallas con una representación de servicio relativamente robusta bajo el modelo \(P\), efecto explicado por la redundancia entre rutas.}
\end{enumerate}

\section{Respuestas a las Preguntas de Investigación}

Retomando las interrogantes planteadas en el Capítulo 1, los modelos construidos nos permiten ofrecer respuestas cuantitativas:

\begin{description}
    \item[¿Cómo está tejida la red?]
    La red presenta una naturaleza dual. En su nivel físico (tejido de infraestructura), es una red dispersa y jerárquica, dependiente de pocas arterias principales. Sin embargo, en su nivel funcional modelado (tejido de servicio proyectado), se comporta como una red compleja de mundo pequeño, facilitando conexiones topológicas entre nodos distantes gracias a la integración de rutas, aunque a un costo elevado de redundancia.

    \item[¿Existen puntos críticos cuya falla colapsaría el servicio?]
    Sí existen puntos críticos en el modelo, aunque el alcance de la conclusión depende de la representación. El análisis de robustez distinguió entre fallos aleatorios y ataques dirigidos. La degradación más severa aparece cuando se remueven supernodos de alta conectividad, especialmente en la infraestructura física (Espacio $L$). En el Espacio $P$, la eliminación de esos nodos no produce colapso global de la proyección, pero este resultado describe conectividad potencial por rutas compartidas, no continuidad operacional completa.

    \item[¿Se agrupan los supernodos en comunidades funcionales?]
    La detección de comunidades en el Espacio $P$ sugiere que la organización funcional del transporte obedece más a corredores de movilidad y zonas de atracción económica que a divisiones administrativas municipales estrictas. Los \textit{cliques} locales observados indican zonas de alta integración interna, pero conectadas laxamente entre sí por las rutas troncales.
\end{description}

\section{Matriz de Trazabilidad de Preguntas de Investigación}

La Tabla~\ref{tab:trazabilidad_preguntas} resume la trazabilidad entre preguntas, métricas observadas, resultados, implicaciones y límites del estudio.

\begin{table}[h!]
    \centering
    \scriptsize
    \renewcommand{\arraystretch}{1.2}
    \setlength{\tabcolsep}{3pt}
    \begin{tabular}{@{}p{0.15\textwidth} p{0.20\textwidth} p{0.18\textwidth} p{0.18\textwidth} p{0.15\textwidth}@{}}
        \toprule
        \textbf{Pregunta} & \textbf{Métrica / evidencia} & \textbf{Resultado clave} & \textbf{Implicación} & \textbf{Límite} \\
        \midrule
        ¿Cómo está tejida la red? &
        L-espacio: densidad $0.000486$, clustering $0.1088$, diámetro $77$. P-espacio: densidad $0.0436$, clustering $0.4085$, diámetro $6$, camino medio $2.41$. &
        Estructura dual: infraestructura dispersa y frágil frente a proyección de servicio densa y de mundo pequeño. &
        La experiencia de conectividad del usuario depende de la superposición de rutas más que de la malla física. &
        El modelo es estático (GTFS programado) y no incorpora variación temporal intradía de la oferta. \\
        \midrule
        ¿Existen puntos críticos cuya falla colapsaría el servicio? &
        Ataque dirigido (50 supernodos de mayor grado), contrastado con fallos aleatorios como línea base. L-espacio: ataque dirigido LCC $97.45\%$, eficiencia $81.50\%$; fallos aleatorios LCC medio $98.23\%$, eficiencia media $97.17\%$. P-espacio: ataque dirigido LCC $98.74\%$, eficiencia $97.36\%$; fallos aleatorios LCC medio $98.80\%$, eficiencia media $96.93\%$. &
        La criticidad es alta en infraestructura y moderada en la proyección $P$; el modelo no colapsa globalmente en $P$, pero sí se degrada en $L$. &
        Conviene priorizar protección y refuerzo de hubs estructurales (p. ej., centro y corredores articuladores). &
        El $P$-espacio por cliques mide conectividad potencial; no se modelan fallos en cascada, frecuencias ni redistribución dinámica de flujo ante interrupciones. \\
        \midrule
        ¿Se agrupan los supernodos en comunidades funcionales? &
        Louvain: 40 comunidades en L-espacio; 11 en P-espacio (modularidad $0.54$), con 67\% de supernodos en las cinco mayores. &
        Sí, existen comunidades funcionales policéntricas asociadas a corredores y zonas de actividad, no estrictamente a límites administrativos. &
        El diseño y gestión puede abordarse por subsistemas funcionales, no sólo por municipio. &
        La partición depende del algoritmo y de los pesos elegidos; otras configuraciones pueden variar los cortes. \\
        \bottomrule
    \end{tabular}
    \caption{Matriz de trazabilidad entre preguntas de investigación, evidencia cuantitativa, hallazgos e implicaciones.}
    \label{tab:trazabilidad_preguntas}
\end{table}

\section{Limitaciones y Agenda de Investigación Futura}

A pesar de la robustez del análisis topológico estático presentado, la abstracción realizada impone límites que abren líneas prometedoras para futura investigación en matemáticas aplicadas al transporte:

\subsection{De Grafos Estáticos a Grafos Temporales ($G_t$)}
Este trabajo modeló la red como una instantánea estática basada en la programación GTFS. Una extensión natural es la modelación mediante \textbf{grafos temporales} o redes variantes en el tiempo, donde las aristas $E(t)$ existen solo si hay un servicio activo en el instante $t$.
\begin{itemize}
    \item \textbf{Investigación Propuesta:} Analizar la fluctuación de la \textit{Eficiencia Global} $Ef(t)$ a lo largo del día para detectar pérdidas de conectividad funcional en horas valle, donde la redundancia representada en el Espacio $P$ podría reducirse y acercar la operación efectiva a la fragilidad del Espacio $L$.
\end{itemize}

\subsection{Análisis Espectral de Grafos}
El presente trabajo se basó principalmente en métricas topológicas clásicas (centralidad, coeficiente de clustering, etc.). La Teoría Espectral de Grafos ofrece una perspectiva complementaria para entender la estructura de la red a través de los autovalores y autovectores de sus matrices asociadas (ej. Laplaciana o de adyacencia).
\begin{itemize}
    \item \textbf{Investigación Propuesta:} Analizar el espectro de las matrices Laplaciana y de Adyacencia de los grafos L-espacio y P-espacio. Esto permitiría cuantificar propiedades globales como la conectividad algebraica, detectar estructuras de partición de manera algorítmica (cortes óptimos) y caracterizar la difusión de información o congestión a través de la red, ofreciendo una visión más profunda sobre la robustez y la dinámica del flujo de pasajeros.
\end{itemize}

\subsection{Matrices de Demanda (Grafos Ponderados por Flujo)}
Nuestro modelo actual asume que todas las aristas tienen igual capacidad o importancia operativa, ignorando la demanda real de pasajeros.
\begin{itemize}
    \item \textbf{Investigación Propuesta:} Integrar matrices Origen-Destino (OD) estimadas para ponderar los nodos no solo por su grado ($k$), sino por su \textbf{carga (load centrality)}. Esto permitiría pasar de un análisis de \textit{vulnerabilidad topológica} a uno de \textit{vulnerabilidad operativa}, identificando nodos que, aunque no sean geométricamente centrales, podrían colapsar el sistema por saturación de usuarios.
\end{itemize}

\subsection{Modelado de Fallos en Cascada (Dinámica No Lineal)}
El análisis de robustez presentado utilizó la eliminación estática de nodos (percolación inversa). Sin embargo, en redes reales, el fallo de un nodo redistribuye el flujo hacia sus vecinos.
\begin{itemize}
    \item \textbf{Investigación Propuesta:} Implementar modelos de \textbf{fallos en cascada} (como el modelo de Motter-Lai) sobre la red de Guadalajara. El objetivo sería determinar la capacidad crítica de los nodos de transbordo necesaria para evitar el colapso sistémico ante perturbaciones locales.
\end{itemize}

\subsection{Optimización Topológica mediante Recableado}
Dado que se identificó un rodeo geométrico moderado (Detour Index \(=1.20\)) y dependencia del centro, existe una oportunidad clara para la optimización matemática.
\begin{itemize}
    \item \textbf{Investigación Propuesta:} Utilizar algoritmos de recableado de aristas (\textit{edge rewiring}) que maximicen la eficiencia global del Espacio $L$ sujeto a restricciones de costo (manteniendo constante la longitud total de aristas). Matemáticamente, esto implica encontrar la topología óptima que minimice el diámetro de la red sin requerir un aumento en la infraestructura física existente.
\end{itemize}

\subsection{Recomendaciones accionables para la planeaci\'on}
Con base en los resultados estructurales y de robustez del Cap\'itulo 4 (Figuras~\ref{fig:centralidad_b_l_space}, \ref{fig:centralidad_b_p_space}, \ref{fig:comunidades_p_space}, \ref{fig:robustez_l_space}, \ref{fig:robustez_p_space}; Tablas~\ref{tab:intervenciones_topologicas_l_space}, \ref{tab:sensibilidad_l_space} y \ref{tab:rangos_consolidados_p3_p4}) y en las decisiones de construcci\'on de datos del Cap\'itulo 3 (Tablas~\ref{tab:construccion_grafo} y \ref{tab:consolidacion_grafo_lspace}), se proponen tres l\'ineas de acci\'on priorizadas:

\begin{enumerate}
    \item \textbf{Prioridad 1: implementar primero refuerzo perif\'erico y protecci\'on de hubs estructurales en el Espacio $L$.}

    \textit{Impacto esperado:} elevar la resiliencia operativa donde hoy se observa la mayor fragilidad, aprovechando que el escenario de refuerzo perif\'erico reporta mejor retenci\'on de eficiencia bajo ataque (88.52\% frente a 82.50\% en el escenario base) y menor di\'ametro post-ataque (53 frente a 84), mientras que la protecci\'on de hubs mejora el LCC residual (+0.79 pp) seg\'un la Tabla~\ref{tab:intervenciones_topologicas_l_space}.

    \textit{Riesgo principal:} sobreinvertir en enlaces nuevos con baja utilizaci\'on efectiva o transferir cuellos de botella hacia nodos no intervenidos.

    \textit{Dependencia de datos:} actualizaci\'on peri\'odica de \texttt{stop\_times\_cleaned.csv} y validaci\'on de conectividad del grafo (Cap\'itulo 3, Tablas~\ref{tab:construccion_grafo} y \ref{tab:consolidacion_grafo_lspace}), complementada con observaciones operativas de tiempos de recorrido para verificar que la mejora topol\'ogica se traduzca en mejora temporal.

    \item \textbf{Prioridad 2: establecer un programa de continuidad operativa para supernodos de transferencia cr\'itica.}

    \textit{Impacto esperado:} reducir la p\'erdida de desempe\~no cuando fallan nodos de alta centralidad, focalizando mantenimiento, gesti\'on de incidencias y protocolos de contingencia en los supernodos que concentran simult\'aneamente grado e intermediaci\'on (Figuras~\ref{fig:centralidad_b_l_space} y \ref{fig:centralidad_b_p_space}), en coherencia con la brecha observada entre la ca\'ida de eficiencia del Espacio $L$ y del Espacio $P$ (Figuras~\ref{fig:robustez_l_space} y \ref{fig:robustez_p_space}).

    \textit{Riesgo principal:} profundizar la dependencia del hipercentro si la continuidad operativa se dise\~na solo para nodos centrales y no incorpora nodos articuladores perif\'ericos.

    \textit{Dependencia de datos:} monitoreo continuo de incidentes y cumplimiento de servicio por supernodo, junto con recalculo peri\'odico de centralidades para mantener vigente la jerarqu\'ia de criticidad.

    \item \textbf{Prioridad 3: institucionalizar la planeaci\'on por comunidades funcionales, no solo por l\'imites administrativos.}

    \textit{Impacto esperado:} mejorar la coordinaci\'on interzonal de oferta y seguimiento operativo en subsistemas de movilidad, aprovechando la evidencia de 11 comunidades funcionales en el Espacio $P$ (modularidad 0.54; Figura~\ref{fig:comunidades_p_space}) y la concentraci\'on del 67\% de supernodos en cinco macrocomunidades.

    \textit{Riesgo principal:} rigidez de dise\~no institucional si la partici\'on comunitaria se toma como fija, pese a que su configuraci\'on depende del algoritmo y de los pesos utilizados.

    \textit{Dependencia de datos:} rec\'alculo programado de comunidades y verificaci\'on de estabilidad con an\'alisis de sensibilidad (Tablas~\ref{tab:sensibilidad_l_space} y \ref{tab:rangos_consolidados_p3_p4}), adem\'as de integrar informaci\'on OD para validar que la delimitaci\'on funcional coincide con la demanda real. Para hacer reproducible esta verificaci\'on, el repositorio p\'ublico del proyecto conserva los insumos GTFS, los scripts de construcci\'on y an\'alisis, y las salidas tabulares y gr\'aficas que sustentan los valores reportados en el Cap\'itulo~4.
\end{enumerate}
